\section{为什么机器人也想要 Foundation Model}

语言和视觉领域已经展示规模化训练带来的复用能力，机器人却仍常为每个任务单独采集、训练和调参。开场的标题与实验室 montage 把这个落差视觉化：机器人能完成许多漂亮演示，但能力增长仍像一组彼此割裂的项目，而不是一条可持续扩展的统一曲线。

\subsection{从一组机器人 Demo 到一个 Scaling 问题}

标题 slide 把方向命名为 “Towards a Robotics Foundation Model”，其中 `Towards` 比 `Foundation Model` 更重要。第二张 montage 展示多台机械臂在不同厨房、物体和任务上的经验，提示真实世界控制已经有大量局部能力；真正的问题是这些经验能否被一个可扩展系统共同吸收。

\lecturefigure{slide-01-title.jpg}{Towards a Robotics Foundation Model：课堂研究议程而非完成声明}{官方录像恢复幻灯片 V001；Stanford CS25 Lecture 15}

\lecturefigure{slide-02-robot-experience-montage.jpg}{真实机器人经验的多样性：不同厨房、物体、动作与 embodiment}{官方录像恢复幻灯片 V002；Stanford CS25 Lecture 15}

\teachervoice{Ted Xiao 一开始把这场课描述为自己对未来方向的“尝试理解”。他的目标是问机器人能力能否从线性、项目式进步转为更快的复合增长，而不是声称今天已经出现可验证的通用机器人 emergence。}

\subsection{课堂路线：Ingredient、System、Feedback、Data}

Agenda 先讲三类 ingredient，再依次进入 RT-1、SayCan、Inner Monologue 和 DIAL。这个顺序对应一台完整机器人：RT-1 学低层技能，SayCan 组合技能，Inner Monologue 根据反馈重规划，DIAL 扩展离线数据的语言覆盖，最后再把各层拼回系统地图。

\lecturefigure{slide-03-agenda-foundation.jpg}{课程路线：从 foundation-model ingredients 到四个机器人系统案例}{官方录像恢复幻灯片 V003；Stanford CS25 Lecture 15}

\subsection{Emergence 与 Homogenization 的吸引力}

Foundation model 对机器人最有吸引力的两点是 \term{emergent capabilities} 与 \term{homogenization}。前者指规模增加后出现小模型不可见的复杂能力；后者指同一个表示和训练接口能够服务组合式下游任务。slide 4 左侧借用语言模型 emergence plots，右侧用“房间—建筑—世界”示意机器人状态空间的巨大跨度。

\lecturefigure{slide-04-why-foundation-model.jpg}{为什么想要 robotics foundation model：emergence 与 downstream homogenization}{官方录像恢复幻灯片 V004；Stanford CS25 Lecture 15}

\begin{knowledgebox}{读图：这不是机器人 scaling curve}
左下曲线来自语言模型任务，用来说明能力可能在 scale 后显现；右侧椭圆只是机器人环境覆盖的概念图。它们共同提出假设：机器人也许需要跨过某个数据、任务与模型规模，才会出现更强组合泛化。课堂没有给出机器人自己的临界点，因此不能把这张图当成已验证 law。
\end{knowledgebox}

\teachervoice{有学生当场问“机器人 foundation model 是否已经存在”。讲者的回答是还不清楚，并进一步说，机器人领域也许需要先看到自己的 emergent-capability curve，才有资格把系统称为 foundation model。}

\subsection{为什么机器人是最后一块 Domino}

语言、图像、代码和音乐可以从互联网获得巨大静态语料，机器人数据却要求真实设备与世界交互。每条轨迹包含硬件时间、操作员成本、安全风险、物体布置和长时状态变化。Domino slide 因此不是说 robotics 没有进展，而是说它缺少与 web-scale token 同样便宜、可复制的数据生产机制。

\lecturefigure{slide-05-why-not-foundation-model.jpg}{Why not：物理交互使 robotics 成为最难倒下的 foundation-model domino}{官方录像恢复幻灯片 V005；Stanford CS25 Lecture 15}

\begin{warningbox}{“更多数据”不是一句可执行建议}
机器人 episode 的价值取决于新增了什么：新任务、新物体、新背景、新 embodiment、新语言还是同一动作的重复。若只增加容易场景的重复轨迹，成本上升却不一定扩大行为覆盖；后面的 RT-1 曲线会直接验证 task diversity 的重要性。
\end{warningbox}

\subsection{把愿景拆成三种 Ingredient}

与其直接训练一个巨型机器人模型，讲者先把问题拆成三类可操作 ingredient：借鉴 ML scaling 的 architecture / tokenization；利用外部 Internet-scale models 的 language 与 common-sense priors；把机器人学习从昂贵在线试错转向可复用的大型 offline datasets。

\lecturefigure{slide-06-ingredients-overview.jpg}{机器人 foundation-model recipe 的三类原料}{官方录像恢复幻灯片 V006；Stanford CS25 Lecture 15}

\begin{importantbox}{三类 Ingredient 不能互相替代}
高容量架构解决表示与互操作问题，外部 foundation models 提供语义 prior，offline robot data 提供真实动作与接触动力学。只有语言没有执行能力，只有机器人轨迹没有语义覆盖，只有大模型没有实时预算，都会让系统停在某一层。
\end{importantbox}

\subsection{本章小结}

机器人向往 foundation model，是因为 emergence 与 homogenization 能把局部技能变成可复用能力；物理数据的昂贵和异质性又使这条路不能照搬互联网训练。课堂于是把愿景拆为架构、外部 prior 和离线数据三种原料。

